\chapter{Технологическая часть}

%17455

\section{Реализация алгоритмов}
Для реализации алгоритмов использовался язык программирования C++, так как он поддерживает статическую типизацию, имеет возможность отключения сборщика мусора и поддерживаеь нативную многопоточность (библиотека threads~\cite[с.~1217]{lit1}). Отладка происходила в среде разработки <<Visual Studio 2022>>.

\lstinputlisting[firstline=172, lastline=187, style=cpp]{../code/graph.cpp}
Листинг 1 -- Последовательный алгоритм.

\lstinputlisting[firstline=189, lastline=243, style=cpp]{../code/graph.cpp}
Листинг 2 -- Параллельный алгоритм.

\section{Тестирование}

\begin{table}[!htbp]
	\caption{Модульные тесты для алгоритмов.}
	\centering
	
	\begin{tabularx}{\textwidth}{|*{3}{>{\centering\arraybackslash}X|}}
		\hline
		\textbf{Алгоритм} &  \textbf{Вход} & \textbf{Выход}\\\hline
		Последовательный & testgraph.gv & 
		\begin{minipage}{5cm}\scriptsize
			\{ n26 n43 n0 n39 n44 n18 n32 n35 n31 n38 n36 n40 n33 n48 n27 n19 n37 n2 n46 n12 n14 n16 \}, 
			\{ n5 n17 n45 n30 n20 \}, \{ n21 n1 n25 \}, \{ n34 n24 n7 \}, \{ n9 n6 \}, \{ n4 n29 \}, \{ n15 n49 \}
		\end{minipage}
		 \\\hline
		Параллельный & testgraph.gv & 
		\begin{minipage}{5cm}\scriptsize
			\{ n26 n43 n0 n39 n44 n18 n32 n35 n31 n38 n36 n40 n33 n48 n27 n19 n37 n2 n46 n12 n14 n16 \}, 
			\{ n5 n17 n45 n30 n20 \}, \{ n25 n1 n21 \}, \{ n34 n24 n7 \}, \{ n4 n29 \}, \{ n9 n6 \}, \{ n15 n49 \} 
		\end{minipage}
		\\\hline
	\end{tabularx}
\end{table}

\section*{Вывод}

В данном разделе мной были реализованы последовательный и параллельный алгоритмы и проведено их успешное тестирование.

\clearpage
